% Generated draft: requires mathematical and source review.
\documentclass[11pt]{article}
\usepackage[margin=25mm]{geometry}
\usepackage{fontspec}
\setmainfont{Latin Modern Roman}
\usepackage{amsmath,amssymb,mathtools,unicode-math}
\setmathfont{Latin Modern Math}
\usepackage{enumitem,needspace,xurl,hyperref}
\setlength{\parindent}{0pt}
\setlength{\parskip}{5pt}
\setlength{\emergencystretch}{3em}
\newcommand{\N}{\mathbb{N}}
\newcommand{\Z}{\mathbb{Z}}
\newcommand{\Q}{\mathbb{Q}}
\newcommand{\R}{\mathbb{R}}
\newcommand{\C}{\mathbb{C}}
\newcommand{\F}{\mathbb{F}}
\newcommand{\A}{\mathbb{A}}
\newcommand{\PP}{\mathbb{P}}
\newcommand{\E}{\mathbb{E}}
\newcommand{\eps}{\varepsilon}
\newcommand{\dd}{\,\mathrm{d}}
\begin{document}
\section*{UnsolvedMath ID: 30004301}
\textbf{Source problem code:} catalogue entry
\subsection{Definitions and mathematical statement}
Let \(X\) be a set and write \(r(x,y)=(\lambda_x(y),\rho_y(x))\) for a map \(r\colon X\times X\to X\times X\). The braid-form Yang--Baxter equation is
\[
(r\times\operatorname{id}_X)(\operatorname{id}_X\times r)(r\times\operatorname{id}_X)
=
(\operatorname{id}_X\times r)(r\times\operatorname{id}_X)(\operatorname{id}_X\times r)
\quad\text{on }X^3.
\]
The structure monoid has the presentation
\[
M(X,r)=\left\langle X\ \middle|\ xy=uv\text{ whenever }r(x,y)=(u,v)\right\rangle.
\]
For a specified commutative ring \(k\) with identity, its monoid algebra \(kM(X,r)\) consists of finite formal \(k\)-linear combinations of elements of \(M(X,r)\), with multiplication extending the monoid multiplication \(k\)-bilinearly. This gives a conventional meaning to the displayed structure algebra; the supplied record does not specify whether the intended coefficient domain is a field or a more general ring. It also does not specify whether the term Yang--Baxter solution requires \(r\) to be bijective. No finiteness or nondegeneracy hypothesis is explicitly imposed in the general question. A solution is called left-nondegenerate when every map \(\lambda_x\colon X\to X\) is bijective.

\textbf{Requested classification, subject to unresolved source qualifications.} For the intended class of set-theoretic Yang--Baxter solutions \((X,r)\), and for the intended associated \(1\)-cocycle \(\pi\), the record requests:
\begin{enumerate}
\item necessary and sufficient conditions for \(\pi\) to be injective;
\item a description of the algebraic structure of \(kM(X,r)\) for solutions satisfying those conditions;
\item a description of the algebraic structure of \(kM(X,r)\) for arbitrary solutions in the intended class.
\end{enumerate}
The record does not define \(\pi\), its domain and codomain, the action entering its cocycle identity, or the cocycle convention. Thus the first request cannot yet be expressed as a predicate on fully specified data. Moreover, both supplied formulations call \(\pi\) an \emph{injective} \(1\)-cocycle before asking when it is injective. Read literally, injectivity is already an assumption and the first request is tautological. The three requests above record the apparent classification target without resolving that conflict. The phrase ``algebraic structure'' also does not identify particular invariants or an equivalence relation for the requested classification.

The supplied dated assessment labels the problem partially solved and states that, for finite solutions, the associated \(1\)-cocycle is bijective exactly in the left-nondegenerate case, that substantial parts of the structure-algebra question have been determined, and that the three classifications remain open in full generality. These are supplied curation claims, not independently verified facts or additional hypotheses; the bijectivity claim does not itself settle the requested injectivity criterion.
\subsection{Short English statement}
The problem asks when an associated cocycle is injective and how to describe the structure algebra in that case and for general Yang--Baxter solutions. A complete precise formulation requires the missing cocycle data, the intended hypotheses on the solutions and coefficients, and the meaning of the requested algebraic classification. The source also assumes injectivity while asking to characterize it. Its partially solved label and literature assessment are supplied dated curation.
\subsection{Sources}
\begin{itemize}
\item Local registry: \texttt{data/unsolvedmath/problems.json}, supplied record with permanent UnsolvedMath ID \(30004301\), title ``Structure Algebras of Yang--Baxter Solutions,'' exact source problem code \texttt{catalogue entry}. The record belongs to the Oberwolfach Reports Open Problems collection. No individual proposer is supplied; the proposed year is recorded as 2019.
\item Supplied report attribution: \emph{Mini-Workshop: Algebraic Tools for Solving the Yang--Baxter Equation} (2020), Oberwolfach Reports, pp.\ 3207--3242, DOI \texttt{10.4171/owr/2019/51}. \url{https://doi.org/10.4171/owr/2019/51}. The report citation is supplied provenance; no problem-specific page or excerpt is provided.
\item Ced\'o--Jespers--Verwimp, ``structure monoids'' (supplied reference label): \url{https://arxiv.org/abs/1912.09710}.
\item Colazzo--Jespers--Kubat--Van Antwerpen, ``structure algebras'' (supplied reference label): \url{https://arxiv.org/abs/2305.06023}.
\item Colazzo--Jespers--Kubat--Van Antwerpen, ``2026 structure algebras'' (supplied reference label): \url{https://doi.org/10.1016/j.jalgebra.2026.04.044}.
\item Literature assessment date: \texttt{2026-08-22}. The partially solved status, the finite-solution bijectivity characterization, the reported progress on structure algebras, and the assertion that all three classifications remain open in full generality are supplied dated curation, not independently verified here. The supplied background expressly requests independent verification of time-sensitive status and cited claims.
\end{itemize}
\end{document}
